\section{Introduction}\label{sec:intro}

Low-light video enhancement (LLVE) is ill-posed: the same dark clip admits many plausible well-lit versions, and supervised models resolve the ambiguity with statistics learned from one paired dataset. They do well on held-out clips of that dataset and degrade on other sensors and scenes~\citep{bvi_rlv}. A natural remedy is a prior that does not come from the paired data. Large frozen encoders are an obvious source, and among them video world models are attractive for video: V-JEPA~\citep{vjepa,vjepa2} is trained to predict the latent representation of masked spatio-temporal regions, so it models how scenes evolve rather than how pixels look.

We ask what such a prior actually buys for LLVE, and answer with controls that the conditioning literature rarely reports. First, we measure illumination invariance of six frozen encoders on paired low-light/well-lit clips (Sec.~\ref{sec:prelim}). V-JEPA is far more invariant than CLIP and SigLIP\,2 but statistically tied with DINOv2/v3 and VideoMAE; what sets it apart is its predictor, which forecasts future latent tokens of dark clips far better than copying the last frame, while VideoMAE's pixel decoder does worse than copying. Second, we build \textbf{Seer}, which conditions a trainable Restormer-style pixel branch on frozen V-JEPA tokens through residual cross-attention, a learnable cross-layer router, and a counterfactual predictor that reuses V-JEPA's own predictor to map the tokens of the dark clip towards those of its well-lit counterpart (Sec.~\ref{sec:method}). Third, we evaluate on six levels from in-domain video to unseen image datasets, with a protocol that fixes a frame-pairing error of earlier evaluations, reports which test sequences each released checkpoint was trained on, and measures temporal consistency with a motion-compensated warping error (Sec.~\ref{sec:exp:protocol}).

The controlled study (Sec.~\ref{sec:exp:control}) keeps initialisation, schedule and seeds fixed and changes only the tokens the pixel branch reads. Pretrained V-JEPA tokens improve SDSD-indoor by $\R{deltanovjepa@L0@psnr_rgb}$\,dB, lower LPIPS on both SDSD splits and make the output more temporally consistent, and a randomly initialised encoder of the same size does not reproduce the gain. On unseen datasets, however, the same tokens do not help on their own. % CROSS-DOMAIN SENTENCE PENDING (GT-mean / few-shot adaptation results, Sec. 5.3)

\paragraph{Contributions.}
(i)~A six-encoder study showing that illumination invariance is common to self-supervised encoders, whereas latent prediction under darkness is specific to the video world model.
(ii)~Seer, an LLVE model that reads a frozen V-JEPA through cross-attention, a cross-layer router and a counterfactual predictor, and is the best of all evaluated methods on the SDSD levels in fidelity, perceptual quality and temporal consistency.
(iii)~A controlled attribution of these gains to the pretrained prior, with random-encoder, encoder-swap and component controls, and an analysis of why they do not carry over to unseen datasets.
(iv)~An evaluation protocol with corrected frame pairing, per-checkpoint training-set membership, a retrained baseline on the same split and a motion-compensated temporal metric; all numbers are generated from released per-sequence results.
